\documentclass[10pt,a4paper]{amsart}
\usepackage[margin=19mm]{geometry}
\usepackage{amsmath,amssymb,mathtools}
\usepackage[scr=rsfs,cal=euler]{mathalfa}
\usepackage{xfrac}
\usepackage[unicode,hidelinks,bookmarks=false]{hyperref}
\newcommand{\ie}{i.e.\ }
\newcommand{\cf}{cf.\ }
\newcommand{\resp}{respectively\ }
\newcommand{\Subsection}{\subsection}
\providecommand{\nobreakdash}{\nobreak\hbox{-}}
\setlength{\emergencystretch}{2em}
\multlinegap0pt
\usepackage{bm}
\usepackage{bbm}
\usepackage{dsfont}
\usepackage{xspace}
\usepackage{cancel}
\usepackage{mathtools}
\usepackage{cleveref}
\usepackage{amsthm}
\makeatletter
\@ifundefined{theorem}{\newtheorem{theorem}{Theorem}}{}
\@ifundefined{corollary}{\newtheorem{corollary}[theorem]{Corollary}}{}
\makeatother
\newcommand{\trans  }{^\top}
\title[The Strong Singular Value Property for Matrices]{The Strong Singular Value Property for Matrices}
\author{Caleb Cheung, Bryan Shader}
\date{Source: arXiv:2507.08313v1 (CC BY 4.0)}
\begin{document}
\maketitle

\noindent\emph{Source and licence.} The Strong Singular Value Property for Matrices. Version arXiv:2507.08313v1 (CC BY 4.0), distributed under the Creative Commons Attribution 4.0 International licence, as recorded in the arXiv licence metadata for this exact version and shown on the abstract page https://arxiv.org/abs/2507.08313v1. Full rights notice: licenses/arxiv-2507.08313-CC-BY-4.0.txt.\par\smallskip

\noindent\textit{Editorial scope.} Counted unit: the Corollary whose source label is \texttt{None} in the article (source0.tex lines 1314--1317, source label \texttt{None}), delivered together with the article's own complete original proof of it (source0.tex lines 1319--1506, one \texttt{proof} environment of 188 lines). No mathematics is written, repaired or re-derived: statement and proof are the article's own text line for line. Required context retained uncounted: ortho (lines 192--200); Path (lines 294--297); distinct (lines 1194--1201). Every definition, display and label of those lines is retained unchanged. Omitted from this package and not redistributed: 1--191, 201--293, 298--1193, 1202--1313, 1318--1318, 1507--2018. Nothing inside a retained excerpt is omitted or altered: every retained body is delivered exactly as the article's lines are written, so no omission receipt of any kind accompanies this package. Citations: the retained excerpts carry no citation command at all, so no bibliography entry of the article is transcribed and no reference is redistributed. Reference adaptation: none was needed and none was made; every reference inside the delivered unit resolves inside it, so no unresolved reference or citation is produced. Printed slips kept literal: no printed word, symbol, hypothesis or number of the counted statement or of its proof was altered. Layout and class: the article's own class is \texttt{elsarticle} with packages including graphicx, dsfont, amsmath, amssymb, amsthm, xcolor, mathabx, float, multirow, hyperref, amsbsy, algorithmic, algorithm2e, cleveref; the wrapper is the lane's pinned \texttt{amsart} editorial wrapper at 10pt on A4, which restates the theorem-like environments the retained text needs and adds the article's own macro definitions that the retained text needs (\texttt{\textbackslash trans}), with extra wrapper packages (bm, bbm, dsfont, xspace, cancel, mathtools, cleveref). The delivered numbering is the unit's own. Assets: no retained span contains a figure environment, a graphics-inclusion command, a PDF-page inclusion, a table or a TikZ picture; no image file is copied into this package, the package declares no asset dependency and no image file is redistributed with it. Licence: version arXiv:2507.08313v1 of this article is distributed under the Creative Commons Attribution 4.0 International licence, as recorded in the arXiv licence metadata for this exact version and shown on the abstract page https://arxiv.org/abs/2507.08313v1. A full rights notice is delivered at licenses/arxiv-2507.08313-CC-BY-4.0.txt. Characters: every retained excerpt is pure ASCII, so the delivered engine, XeLaTeX with its default Latin Modern text font, reports no missing character.
\begin{theorem}
\label{ortho}
Let $P$ be an $m\times n$ pattern with $m\leq n$.
Then there exists a matrix with pattern $P$ having 
all singular values equal and nonzero if and only if 
there is a row-orthonormal matrix with pattern $P$.
Moreover, in this case,   each list of $m$ positive real numbers is 
the list of singular values of a matrix with pattern $P$.
\end{theorem}

\begin{theorem}
\label{Path}
Let $P$ be an $n\times (n+1)$ pattern whose underlying bigraph is a path on $2n+1$ vertices.  Then a list $\Sigma$ of $n$ nonnegative real numbers forms the singular values of a matrix with pattern $P$ if and only if the elements are distinct and nonzero.
\end{theorem}

\begin{corollary}[\bf Distinct Singular Values Theorem]
\phantom{ }\\ 
\label{distinct}
Let $\Sigma$ be a set of $m$ distinct positive real numbers 
and $P$ be an $m\times n$ pattern. 
Then there exists an $m\times n$ matrix $A$ with pattern $P$
and singular values $\Sigma$ if and only if $P$ has term-rank $m$.
\end{corollary}

\begin{corollary}
There exists a matrix with pattern $C_6$ and singular values $\sigma_1\geq \sigma_2 \geq \sigma_3 \geq 0$ if and only 
if $\sigma_2>0$ and $\sigma_1 \neq  \sigma_3$.
\end{corollary}

\begin{proof}
As $C_6$ is not the pattern of a row-orthogonl matrix,
by \Cref{ortho} there is no matrix with pattern $C_6$
having all singular values equal. As every matrix with pattern $C_6$ has rank at least two, there is no matrix with pattern $C_6$ having $0$ as a singular value of multiplicity 2.  Hence the conditions on the $\sigma$ are necessary. 

\bigskip\noindent
We prove sufficiency by cases.

\bigskip\noindent
{\bf Case 1.} $ \sigma_1>\sigma_2> \sigma_3>0$. 
\\ 
By \Cref{distinct}, there exists a matrix with pattern $C_6$ and these singular values.

\bigskip\noindent
{\bf Case 2.} $ \sigma_1>\sigma_2> \sigma_3=0$. 
\\ 
There exist positive reals $a, b, c, d$ such that 
 $a^2 \neq c^2$, $a^2+ b^2=\sigma_1^2$, $c^2+d^2=\sigma_2^2$, and  
the matrix 
\[ 
N= \left[ \begin{array}{ccc}
a & b & 0 \\
0 & 0 & c \\
0 & 0 & d 
\end{array}
\right] 
\]
has singular values $\sigma_1$, $\sigma_2$ and $0$. 
To verify that  $N$ has the SSVP property with respect 
to the $C_6$,
let 
\[ 
X=\left[\begin{array}{ccc}
0& 0 & x\\ 
y & 0 & 0\\
0 & z & 0
\end{array} 
\right],
\]
and assume that both $N\trans X$ and $X N\trans$ are symmetric. 
Then 
\begin{eqnarray*}
ax&=&cy\\ 
bx&=&dz\\ 
cx&=&ay\\
bz&=&dx.
\end{eqnarray*}
The first and third of these equations imply that $a^2xy=c^2xy$.
Since $a^2\neq c^2$, $xy=0$. Hence $x=0$ or $y=0$. In either case, 
the first of the equations imply that both $x$ and $y$ equal $0$. 
Now the second equation implies that $z=0$.  Hence, $X=O$, and 
$N$ has the SSVP with respect to the desired pattern. 


Consider the skew-symmetric matrices
\[
K=\left[ 
\begin{array}{ccc}
0 & (d^2-b^2)c & (a^2-c^2)d \\
(b^2-d^2)c & 0 & 0 \\
(c^2-a^2)d & 0 &0 
\end{array} \right] \mbox{ and } 
\]
\[   
L= 
\left[
\begin{array}{ccc}
            0    &          0  & (b^2 - d^2)a\\
             0   &0 & (c^2 - a^2)b\\
(d^2 - b^2)a & (a^2 - c^2)b &              0
\end{array} \right]. 
\]
One can verify that 
\[ 
KN+NL = 
\left[
\begin{array}{ccc}
0 & 0 & 0\\
0 &  (a^2 + b^2 - c^2 - d^2)bc &  0 \\
-(a^2 + b^2 - c^2 - d^2)ad &  0 & 0
\end{array} 
\right]
\] 
and $KN+NL\in \mbox{Tan}_N^{\Sigma}$.
Hence the $(2,2)$ and $(3,1)$ entries of $KN+NL$ are positive and negative respectively.
By the Matrix Liberation Theorem, there exist a matrix whose  pattern is
$C_6$
having singular values $\sigma_1$, $\sigma_2$ and $0$. 

\bigskip\noindent
{\bf Case 3.} $ \sigma_1=\sigma_2> \sigma_3=0$. 
\\ 
Note that  
\[ 
\frac{1}{\sqrt{3}}\left[ 
\begin{array}{rrr}
1 & 1 & 0 \\
0 & 1 & 1 \\
-1 & 0 & 1
\end{array} 
\right]
\]
has pattern $C_6$, and singular values $\sigma_1$, $\sigma_2$ and $0$.

\bigskip\noindent
{\bf Case 4.} $\sigma_1>\sigma_3>0$, and either 
$\sigma_1=\sigma_2$ or $\sigma_2=\sigma_3$.

Without loss of generality we may assume the desired 
singular values are 
 are $\sigma$, $1$ and $1$,
where $\sigma$ is a positive real number not equal to $1$. 
By \Cref{Path}, there exist positive reals $a$, $b$, $c$ such that 
the matrix 
\[ 
N=\left[ 
\begin{array}{cc}
a& b 
\\
0 & c
\end{array}
\right] 
\] 
has singular values $\sigma$ and $1$. Hence the matrix 
\[ 
M=\left[
\begin{array}{ccc}
a & b & 0 
\\
0 & c & 0 \\
0 & 0 & 1 
\end{array}
\right] 
\]
has singular values $1$, $1$ and $\sigma$.

We can verify that $M$ has the SSVP relative to 
pattern 
\[
P=\left[ \begin{array}{ccc}
1 & 1& 0 \\
0 & 1 & 1 \\
1 & 0 & 1
\end{array} 
\right] 
\]
 by considering a matrix $X$ of the 
form
\[ 
X= \left[ \begin{array}{ccc} 0 & 0 & x \\ y & 0 & 0 \\ 0 & z & 0 
\end{array} \right] 
\]
such that $M^TX$ and $XM^T$ are symmetric.
The symmetry of $M^TX$ is equivalent to $cy=0$, $ax=0$ and $bx=z$.
Since, each of $a$, $b$, and $c$ is nonzero, each of $x$, $y$ and $z$ are zero.  Thus, $X=O$.

Now let 
\[
K= 
\left[ \begin{array}{rrr}
0 & 0 & -1 \\
0 & 0 & 0 \\
1 & 0 & 0 \end{array} 
\right] 
\mbox{ and } L= \left[ \begin{array}{ccc} 
0 & 0 & (1-b^2)/a 
\\
0 & 0 & b \\
(b^2-1)/a & -b & 0 
\end{array} \right].
\]
Then 
\[ 
KM+ML= \left[ 
\begin{array}{ccc}
0 & 0 & 0 \\
0 & 0 & bc\\
(a^2 + b^2 -1)/a & 
0 & 0 
\end{array} 
\right],
\]
and $KM+ML$ is in $\mbox{Tan}^{\Sigma}_A$.


Thus, by the Matrix Liberation Theorem,
there exists a matrix $B$ with singular values $1$, $1$ and $\sigma$ whose pattern is that of $A$ in the direction of $KN+NL$. 
\end{proof}

\end{document}
